\section*{סיכום הרצאה: חלקיקים זהים ותורת הפיזור}

\subsection*{1. עקרון האיסור של פאולי ואופרטור האנטי-סימטריזציה}

בהרצאה זו המשכנו את הדיון במערכות של חלקיקים זהים, תוך התמקדות בפרמיונים (\LR{Fermions}). פונקציית הגל של פרמיונים חייבת להיות אנטי-סימטרית לחלוטין תחת החלפה של כל זוג חלקיקים.

\begin{definitionbox}{אופרטור האנטי-סימטריזציה (\LR{Antisymmetrizer})}
עבור מערכת של $n$ חלקיקים, נגדיר את האופרטור $\hat{\mathcal{A}}$ באופן הבא:
$$\hat{\mathcal{A}} = \frac{1}{n!} \sum_{P \in S_n} \text{sgn}(P) \hat{P}$$
כאשר:
\begin{itemize}
    \item $S_n$ היא חבורת התמורות (\LR{Permutation Group}) של $n$ איברים.
    \item $\hat{P}$ הוא אופרטור פרמוטציה.
    \item $\text{sgn}(P)$ הוא סימן הפרמוטציה: $1$ עבור פרמוטציה זוגית ו-$-1$ עבור אי-זוגית.
\end{itemize}
\end{definitionbox}

\subsubsection*{תכונות האופרטור}
עבור אופרטור החלפה של זוג חלקיקים $\hat{P}_{ij}$, מתקיים:
$$\hat{P}_{ij} \hat{\mathcal{A}} = -\hat{\mathcal{A}}$$
מכאן נובע כי אם נפעיל את האופרטור על מצב שבו שני חלקיקים נמצאים באותו מצב קוונטי (למשל מצב $|C\rangle$ עבור חלקיקים 3 ו-4), נקבל:
$$\hat{\mathcal{A}} |a, b, C, C, \dots\rangle = 0$$
זהו ניסוח מתמטי של \textbf{עקרון האיסור של פאולי} (\LR{Pauli Exclusion Principle}): לא ניתן לאכלס שני פרמיונים זהים באותו מצב קוונטי.

\begin{examplebox}{המערכה המחזורית (\LR{The Periodic Table})}
ניתן להסביר את מבנה הקליפות באטומים באמצעות עקרון פאולי. עבור רמות אנרגיה של אטום דמוי-מימן:
\begin{itemize}
    \item ברמה $n=1$, ישנו ניוון של 2 (בשל הספין). לכן הליום ($Z=2$) סוגר קליפה והוא גז אציל (\LR{Noble Gas}).
    \item ברמה $n=2$, ישנו ניוון של 8. אטום הניאון ($Z=10$) סוגר קליפה נוספת.
    \item אטומים עם אלקטרון בודד מעל קליפה סגורה (כמו ליתיום) מתנהגים כמתכות אלקליות (\LR{Alkali Metals}).
\end{itemize}
\end{examplebox}

\subsection*{2. דטרמיננטת סלייטר (\LR{Slater Determinant})}

כדי לבנות פונקציית גל אנטי-סימטרית מתוך אוסף של מצבים חד-חלקיקיים $\phi_a, \phi_b, \dots$, נשתמש בדטרמיננטה:

\begin{resultbox}{דטרמיננטת סלייטר}
פונקציית הגל המנורמלת של $n$ פרמיונים במצבים $\alpha_1, \dots, \alpha_n$ ניתנת לכתיבה כ:
$$\Psi(x_1, \dots, x_n) = \frac{1}{\sqrt{n!}} \begin{vmatrix} 
\phi_{\alpha_1}(x_1) & \phi_{\alpha_1}(x_2) & \dots & \phi_{\alpha_1}(x_n) \\
\phi_{\alpha_2}(x_1) & \phi_{\alpha_2}(x_2) & \dots & \phi_{\alpha_2}(x_n) \\
\vdots & \vdots & \ddots & \vdots \\
\phi_{\alpha_n}(x_1) & \phi_{\alpha_n}(x_2) & \dots & \phi_{\alpha_n}(x_n)
\end{vmatrix}$$
החלפת שני חלקיקים שקולה להחלפת עמודות (המשנה סימן), וקיום שני מצבים זהים שקול לשתי שורות זהות (המאפסות את הדטרמיננטה).
\end{resultbox}

\subsection*{3. חפיפת פונקציות גל והפרדה מרחבית}

עולה השאלה: האם עלינו לעשות אנטי-סימטריזציה לאלקטרון בירושלים ולאלקטרון בבאר שבע?
מבחינה פורמלית התשובה היא כן, אך מבחינה פיזיקלית האפקט זניח כאשר אין חפיפה (\LR{Overlap}) בין תומכי פונקציות הגל.

\begin{theorembox}{הזנחת אנטי-סימטריזציה במרחקים גדולים}
נניח שני מצבים מרחביים מופרדים $J$ ו-$B$. פונקציית הגל האנטי-סימטרית היא:
$$\Psi(1,2) = \frac{1}{\sqrt{2}} \left( \phi_J(1)\phi_B(2) - \phi_B(1)\phi_J(2) \right)$$
במדידה של אופרטור מקומי, איברי ה"חילוף" (\LR{Exchange terms}) כוללים אינטגרלים מהצורה $\int \phi_J^*(x) \hat{O} \phi_B(x) dx$. אם $\phi_J$ ו-$\phi_B$ אינן חופפות מרחבית, אינטגרלים אלו מתאפסים. במצב זה, החלקיקים מתנהגים כניתנים להבחנה (\LR{Distinguishable}).
\end{theorembox}

\subsection*{4. מבוא לתורת הפיזור (\LR{Scattering Theory})}

ניסויי פיזור (כמו ניסוי רתרפורד) הם הכלי העיקרי לחקר מבנה החומר. הבעיה הטיפוסית כוללת קרן חלקיקים הפוגעת במטרה (\LR{Target}) ומתפזרת לעבר גלאים.

\subsubsection*{תיאור פיזיקלי וסקאלות}
\begin{itemize}
    \item המטרה קטנה מאוד ($10^{-10}$ \text{m}) ביחס למרחק לגלאי (~1 \text{m}).
    \item החלקיק הנכנס מתואר כגל מישורי (\LR{Plane Wave}): $\psi_{inc} \approx e^{i\vec{k}\cdot\vec{x}}$.
    \item החלקיק המפוזר מתואר אסימפטוטית כגל כדורי (\LR{Spherical Wave}).
\end{itemize}


\subsubsection*{חתך פעולה (\LR{Cross Section})}
\begin{definitionbox}{חתך פעולה דיפרנציאלי}
$$\frac{d\sigma}{d\Omega} = \frac{\text{קצב אירועים ליחידת זווית מרחבית}}{\text{שטף נכנס}}$$
היחידות של $\sigma$ הן יחידות שטח (מכאן השם "חתך פעולה").
\end{definitionbox}

\subsubsection*{פיתוח מתמטי אסימפטוטי}
עבור פוטנציאל קצר טווח $U(r)$, פונקציית הגל במרחק גדול מהמטרה נראית כך:
$$\psi(\vec{r}) \xrightarrow[r \to \infty]{} \frac{1}{\sqrt{V}} \left[ e^{ikz} + f(\theta, \phi) \frac{e^{ikr}}{r} \right]$$
כאשר $f(\theta, \phi)$ נקראת \textbf{אמפליטודת הפיזור} (\LR{Scattering Amplitude}).

נחשב את שטף ההסתברות (\LR{Probability Flux}) עבור הגל המפוזר $\psi_{scat} = f \frac{e^{ikr}}{r}$:
$$\vec{J} = \frac{\hbar}{m} \text{Im}(\psi^* \nabla \psi) \approx \frac{\hbar k}{m} \frac{|f(\theta)|^2}{r^2} \hat{r}$$
קצב האירועים בגלאי בעל שטח $dA = r^2 d\Omega$ הוא:
$$dN = J_{scat} \cdot r^2 d\Omega = \frac{\hbar k}{m} |f(\theta)|^2 d\Omega$$
השטף הנכנס הוא $J_{inc} = \frac{\hbar k}{m}$. לכן:

\begin{resultbox}{הקשר בין חתך הפעולה לאמפליטודת הפיזור}
$$\frac{d\sigma}{d\Omega} = |f(\theta, \phi)|^2$$
כל המידע על האינטראקציה הפיזיקלית מקודד בתוך האמפליטודה $f(\theta, \phi)$.
\end{resultbox}

\begin{notebox}{קירובים אסימפטוטיים}
בחישוב השטף הרדיאלי, הזנחנו איברים מסדר $1/r^3$ ומעלה, שכן הם דועכים מהר מאוד ואינם תורמים למדידה במרחק גדול (התחום הפרוקסימלי לעומת התחום הרחוק).
\end{notebox}
\newpage

\section{חלקיקים זהים, אנטי-סימטריזציה וחוק האיסור של פאולי}

\subsection{חלקיקים זהים וחבורת הפרמוטציות}

נבחן מערכת של $N$ חלקיקים זהים. חילוף של שני חלקיקים מתואר על-ידי אופרטור פרמוטציה $P\in S_N$, כאשר $S_N$ היא \LR{Symmetric Group} של כל התמורות של $N$ עצמים.

כל פרמוטציה ניתנת לפירוק למכפלה של החלפות זוגיות (טרנספוזיציות). נסמן ב־$\mathrm{sgn}(P)=\pm1$ את \LR{signature} של הפרמוטציה:
\begin{itemize}
    \item $\mathrm{sgn}(P)=+1$ אם מספר ההחלפות זוגי.
    \item $\mathrm{sgn}(P)=-1$ אם מספר ההחלפות אי-זוגי.
\end{itemize}



\subsection{אופרטור האנטי-סימטריזציה}

\begin{definitionbox}{אופרטור האנטי-סימטריזציה}
עבור מערכת של $N$ פרמיונים, האופרטור מוגדר כסכום על כל הפרמוטציות האפשריות:
\[
A \;=\; \frac{1}{N!}\sum_{P\in S_N} \mathrm{sgn}(P)\,P .
\]
זהו אופרטור הטלה (\LR{projection operator}) המבודד את תת-המרחב של המצבים האנטי-סימטריים.
\end{definitionbox}

\begin{resultbox}{תכונת החילוף}
לכל החלפה של זוג חלקיקים $(i,j)$ מתקיים:
\[
P_{ij}\,A \;=\; A\,P_{ij} \;=\; -A .
\]
\end{resultbox}

\begin{notebox}{הוכחת הסימן}
ההוכחה נשענת על כך שכפל של פרמוטציה כללית בטרנספוזיציה (חילוף בודד) מוסיף החלפה אחת לספירה, ולכן הופך את סימן ה-$\mathrm{sgn}$.
\end{notebox}

\subsection{מצבים פרמיוניים כלליים}
כל מצב פרמיוני ניתן לכתיבה בצורה:
\[
\Psi \;=\; \mathcal{N}\, A\,\Phi ,
\]
כאשר $\Phi$ היא פונקציית גל כללית (שאינה בהכרח אנטי-סימטרית), ו-$\mathcal{N}$ קבוע נרמול. אם $A\Phi\neq 0$, מתקבל מצב פרמיוני פיזיקלי לגיטימי.



\section{חוק האיסור של פאולי}

\subsection{הוכחה מתמטית}

נניח כי בפונקציית הגל $\Phi$ שני חלקיקים (למשל $i$ ו-$j$) מאכלסים את אותו מצב חד-חלקיקי בדיוק. אזי חילוף של שני החלקיקים אינו משנה את $\Phi$:
\[
P_{ij}\Phi = \Phi .
\]

מאידך, מהתכונה של אופרטור האנטי-סימטריזציה שהראינו לעיל:
\[
P_{ij}(A\Phi)= -A\Phi .
\]

שילוב שתי המשוואות גורר סתירה אלא אם הפונקציה מתאפסת:
\[
A\Phi = -A\Phi \;\;\Rightarrow\;\; 2A\Phi = 0 \;\;\Rightarrow\;\; \Psi = 0.
\]

\begin{theorembox}{חוק האיסור של פאולי}
לא ייתכנו שני פרמיונים באותו מצב קוונטי. האנטי-סימטריזציה מאפסת כל מצב שבו יש ניוון מלא במצבי החלקיקים.
\end{theorembox}



\section{יישום: מבנה האטום והטבלה המחזורית}

\subsection{רמות האנרגיה באטום המימן}
בהזנחת תיקונים עדינים (\LR{fine structure}), רמות האנרגיה תלויות רק במספר הקוונטי הראשי $n$. הניוון של רמה נתון על-ידי:
\[
g_n = 2n^2,
\]
כאשר הגורם $2$ נובע מהספין של האלקטרון.

\begin{examplebox}{מילוי קליפות באטומים קלים}
\begin{itemize}
    \item \textbf{מימן ($Z=1$):} אלקטרון יחיד ברמת היסוד.
    \item \textbf{הליום ($Z=2$):} שני אלקטרונים באותה רמה מרחבית ($1s$) אך עם ספינים מנוגדים. זוהי קליפה סגורה.
    \item \textbf{ליתיום ($Z=3$):} החלקיק השלישי אינו יכול להיכנס לרמת היסוד (פאולי), ולכן נאלץ לעלות לרמה הבאה ($2s$).
\end{itemize}
\end{examplebox}



\begin{resultbox}{כימיה קוונטית}
סגירת קליפות אלקטרוניות מסבירה את היציבות של הגזים האצילים ואת המחזוריות בתכונות הכימיות של היסודות.
\end{resultbox}



\section{מתי ניתן להזניח אנטי-סימטריזציה?}

נבחן שני אלקטרונים עם פונקציות גל מרחביות $\phi_J$ ו-$\phi_B$ שמרוכזות באזורים מרוחקים (למשל "ירושלים" ו"באר-שבע").

אם החפיפה (\LR{Overlap}) זניחה:
\[
\langle \phi_J | \phi_B \rangle \approx 0 ,
\]
אזי האיברים הצולבים ("איברי ההחלפה") הנובעים מהאנטי-סימטריזציה מתאפסים בחישובי תוחלת והסתברות.

\begin{mainbox}{הבחנה אפקטיבית}
כאשר החפיפה בין פונקציות הגל זניחה, ניתן להתייחס לפרמיונים כאילו היו חלקיקים לא-זהים (מובחנים על ידי מיקומם).
\end{mainbox}



\section{דטרמיננטת סלייטר}

\subsection{מקרה פרטי: שני חלקיקים}
עבור שני פרמיונים במצבים חד-חלקיקיים $A,B$, הפונקציה האנטי-סימטרית היא:
\[
\Psi(1,2)=\frac{1}{\sqrt{2}}
\begin{vmatrix}
\phi_A(1) & \phi_A(2) \\
\phi_B(1) & \phi_B(2)
\end{vmatrix}
=
\frac{1}{\sqrt{2}}\big[\phi_A(1)\phi_B(2)-\phi_A(2)\phi_B(1)\big].
\]

\subsection{הכללה ל-$N$ חלקיקים}

\begin{definitionbox}{דטרמיננטת סלייטר \LR{(Slater Determinant)}}
פונקציית גל פרמיונית כללית של חלקיקים בלתי-תלויים ניתנת לכתיבה כדטרמיננטה:
\[
\Psi(1,\dots,N)=\frac{1}{\sqrt{N!}}
\begin{vmatrix}
\phi_1(1) & \phi_1(2) & \cdots & \phi_1(N) \\
\phi_2(1) & \phi_2(2) & \cdots & \phi_2(N) \\
\vdots & \vdots & \ddots & \vdots \\
\phi_N(1) & \phi_N(2) & \cdots & \phi_N(N)
\end{vmatrix}.
\]
\end{definitionbox}



\begin{notebox}{תכונות הדטרמיננטה}
\begin{itemize}
    \item \textbf{אנטי-סימטריות:} החלפת שני חלקיקים שקולה להחלפת שתי עמודות $\Rightarrow$ שינוי סימן הדטרמיננטה.
    \item \textbf{איסור פאולי:} אם שני חלקיקים באותו מצב, יש שתי שורות זהות $\Rightarrow$ הדטרמיננטה מתאפסת.
\end{itemize}
\end{notebox}



\section{השלכות פיזיקליות רחבות}
עקרון האיסור הוא הבסיס להבנת מגוון תופעות:
\begin{itemize}
    \item יציבות חומר דחוס (מניעת קריסה של החומר).
    \item לחץ הניוון (\LR{Degeneracy Pressure}) בכוכבי נייטרונים וננסים לבנים.
    \item מבנה פסי האנרגיה במתכות ומוליכים למחצה.
\end{itemize}

\begin{summarybox}{סיכום}
סטטיסטיקת פרמי-דיראק, האנטי-סימטריזציה ודטרמיננטות סלייטר מספקות מסגרת מאוחדת להבנת הטבע הקוונטי של החומר, מהרמה האטומית ועד לאובייקטים אסטרופיזיקליים. עקרון פאולי אינו "כוח" במובן הקלאסי, אלא אילוץ גיאומטרי-טופולוגי על פונקציית הגל.
\end{summarybox}

